\documentclass{article}
\usepackage[UTF8]{ctex}
\usepackage{amsmath,amssymb,amsfonts}
\usepackage{enumitem}
\usepackage{geometry}
\geometry{a4paper, margin=1in}
% === 错题本排版优化（不参与题目标识） ===
\setlength{\parindent}{0pt}
\setlength{\parskip}{0.35em plus 0.10em minus 0.08em}
\setlength{\emergencystretch}{3em}
\setlist{itemsep=0.25em, topsep=0.40em, parsep=0pt}
\allowdisplaybreaks[2]
% === 错题本排版优化结束 ===
\begin{document}
\section*{首都师范大学2024年数学分析真题}
\section*{31．首都师范大学2024年数学分析试卷}

\subsection{题目1}
1． \(\lim _{x \rightarrow 0^{+}}(\cos \sqrt{x})^{\frac{1}{x}}\) ．

2． \(\lim _{n \rightarrow \infty} a_{n}=a>0\) ，求 \(\lim _{n \rightarrow \infty} \sqrt[n]{a_{n}}\) ．

3． \(\lim _{x \rightarrow 0} \frac{(1+\tan x)^{\frac{1}{4}}+(1-\sin x)^{\frac{1}{4}}-2}{x^{2}}\) ．

\subsection{题目2}
二．（10 分）讨论 \(f(x, y)=\left(1+\frac{2}{x^{2}}\right)^{\frac{x^{4}}{x^{2}+y^{2}}}\) 在 \((0,0)\) 处累次极限与重极限的存在性，若存在，求其值．

\subsection{题目3}
三．（10 分）证明：\(f(x)=\ln x\) 在 \((1,+\infty)\) 上一致连续，但在 \((0,1)\) 上不一致连续．

\subsection{题目4}
四．（10 分）设 \(f(x)\) 在 \([a, b]\) 上可导，满足 \(f(a)=f(b)=0, f_{+}^{\prime}(a) f_{-}^{\prime}(b)>0\) ，证明：至少存在不同的两点 \(\xi, \eta \in(a, b)\) ，使得 \(f^{\prime}(\xi)=f^{\prime}(\eta)=0\) ．

\subsection{题目5}
五。（12 分）证明：\(f(x, y)=x^{3}-4 x^{2}+2 x y-y^{2}\) 在 \(\mathbb{R}^{2}\) 上有唯一的极值点，且该极值点是极大值点，不是最大值点。

\subsection{题目6}
六．（14 分）求三重积分
\[
\iiint_{\Omega} 2 \sqrt{x^{2}+y^{2}} \mathrm{~d} x \mathrm{~d} y \mathrm{~d} z
\]
其中 \(\Omega\) 为柱面 \(y=\sqrt{2 x-x^{2}}\) 及平面 \(z=0, z=a(a>0), y=0\) 所围成的区域．

\subsection{题目7}
七．（15 分）定义 \(f(x)= \begin{cases}\frac{1}{x} e^{-\frac{1}{x^{2}}}, & x \neq 0 ; \\ 0, & x=0 .\end{cases}\)

（1） （5 分）证明：\(f(x)\) 在点 \(x=0\) 处连续且可导．
（2） （ 5 分）证明：\(f^{\prime}(x)\) 在 \(\mathbb{R}\) 上连续．
（3） （5 分）求 \(f(x)\) 的单调区间，最大值点，最小值点．

\subsection{题目8}
八．（15 分）函数项级数 \(\sum_{n=2}^{\infty} \frac{x e^{-n x}}{\ln n}\) ．

（1） （7 分）求此函数项级数的收敛区间．
（2） （8 分）设 \(a>0\) ，证明此函数项级数在 \([a,+\infty)\) 上一致收敛．

\subsection{题目9}
九．（15 分）设 \(f_{1}(x)\) 在 \([a, b]\) 上可积，\(A\) 为一给定的数，定义
\[
f_{n+1}(x)=A+\int_{a}^{x} f_{n}(t) \mathrm{d} t
\]
其中 \(x \in[a, b], n=1,2, \cdots\) ．
（1）（8 分）证明：函数列 \(\left\{f_{n}(x)\right\}\) 在 \([a, b]\) 上一致收敛。
（2）（7 分）记 \(\left\{f_{n}(x)\right\}\) 的极限函数为 \(f(x)\) ，证明：\(f(x)\) 在 \([a, b]\) 上可微．

\subsection{题目10}
十．（15 分）计算曲线积分
\[
I=\int_{L} x \ln \left(x^{2}+y^{2}-1\right) \mathrm{d} x+y \ln \left(x^{2}+y^{2}-1\right) \mathrm{d} y
\]
其中 \(L\) 是被积函数定义域内从点 \((2,0)\) 到点 \((0,2)\) 的逐段光滑曲线．

\subsection{题目11}
十一．（10 分）设 \(f(x)\) 在 \([a, b]\) 上可积，证明： \(\lim _{\lambda \rightarrow+\infty} \int_{a}^{b} f(x) \sin \lambda x \mathrm{~d} x=0\) ．
\end{document}
